\documentclass[10pt,a4paper]{amsart}
\usepackage{amsmath,amssymb,amsthm,graphicx}
\usepackage{verbatim}
\newenvironment{reference}{\comment}{\endcomment}
%\newenvironment{reference}{}{}
\newenvironment{slogan}{\comment}{\endcomment}
\newenvironment{history}{\comment}{\endcomment}

% For commutative diagrams we use Xy-pic
\usepackage[all]{xy}

% We use 2cell for 2-commutative diagrams.
\xyoption{2cell}
% Theorem environments.
%
\theoremstyle{plain}
\newtheorem{theorem}[subsection]{Theorem}
\newtheorem{proposition}[subsection]{Proposition}
\newtheorem{lemma}[subsection]{Lemma}

\theoremstyle{definition}
\newtheorem{definition}[subsection]{Definition}
\newtheorem{example}[subsection]{Example}
\newtheorem{exercise}[subsection]{Exercise}
\newtheorem{situation}[subsection]{Situation}

\theoremstyle{remark}
\newtheorem{remark}[subsection]{Remark}
\newtheorem{remarks}[subsection]{Remarks}

\numberwithin{equation}{subsection}

% Macros
%
\def\lim{\mathop{\mathrm{lim}}\nolimits}
\def\colim{\mathop{\mathrm{colim}}\nolimits}
\def\Spec{\mathop{\mathrm{Spec}}}
\def\Hom{\mathop{\mathrm{Hom}}\nolimits}
\def\Ext{\mathop{\mathrm{Ext}}\nolimits}
\def\SheafHom{\mathop{\mathcal{H}\!\mathit{om}}\nolimits}
\def\SheafExt{\mathop{\mathcal{E}\!\mathit{xt}}\nolimits}
\def\Sch{\mathit{Sch}}
\def\Mor{\mathop{\mathrm{Mor}}\nolimits}
\def\Ob{\mathop{\mathrm{Ob}}\nolimits}
\def\Sh{\mathop{\mathit{Sh}}\nolimits}
\def\NL{\mathop{N\!L}\nolimits}
\def\CH{\mathop{\mathrm{CH}}\nolimits}
\def\proetale{{pro\text{-}\acute{e}tale}}
\def\etale{{\acute{e}tale}}
\def\QCoh{\mathit{QCoh}}
\def\Ker{\mathop{\mathrm{Ker}}}
\def\Im{\mathop{\mathrm{Im}}}
\def\Coker{\mathop{\mathrm{Coker}}}
\def\Coim{\mathop{\mathrm{Coim}}}

% Boxtimes
%
\DeclareMathSymbol{\boxtimes}{\mathbin}{AMSa}{"02}

%
% Macros for moduli stacks/spaces
%
\def\QCohstack{\mathcal{QC}\!\mathit{oh}}
\def\Cohstack{\mathcal{C}\!\mathit{oh}}
\def\Spacesstack{\mathcal{S}\!\mathit{paces}}
\def\Quotfunctor{\mathrm{Quot}}
\def\Hilbfunctor{\mathrm{Hilb}}
\def\Curvesstack{\mathcal{C}\!\mathit{urves}}
\def\Polarizedstack{\mathcal{P}\!\mathit{olarized}}
\def\Complexesstack{\mathcal{C}\!\mathit{omplexes}}
% \Pic is the operator that assigns to X its picard group, usage \Pic(X)
% \Picardstack_{X/B} denotes the Picard stack of X over B
% \Picardfunctor_{X/B} denotes the Picard functor of X over B
\def\Pic{\mathop{\mathrm{Pic}}\nolimits}
\def\Picardstack{\mathcal{P}\!\mathit{ic}}
\def\Picardfunctor{\mathrm{Pic}}
\def\Deformationcategory{\mathcal{D}\!\mathit{ef}}
\usepackage{cleveref}
\crefname{originallemma}{original lemma}{original lemmas}
\Crefname{originallemma}{Original lemma}{Original lemmas}
\makeatletter
\newlabel{lemma-localized-chern-pre}{{0F9K}{source}{}{sourcepointer-1}{}}
\newlabel{lemma-localized-chern-pre@cref}{{[originallemma][0][]0F9K}{[1][0][]source}{}{}{}}
\makeatother
\crefname{originalreference}{original reference}{original references}
\Crefname{originalreference}{Original reference}{Original references}
\makeatletter
\newlabel{remark-ring-loc-classes}{{0FA0}{source}{}{sourcepointer-2}{}}
\newlabel{remark-ring-loc-classes@cref}{{[originalreference][0][]0FA0}{[1][0][]source}{}{}{}}
\makeatother
\crefname{originallemma}{original lemma}{original lemmas}
\Crefname{originallemma}{Original lemma}{Original lemmas}
\makeatletter
\newlabel{lemma-gysin-at-infty}{{0F9J}{source}{}{sourcepointer-3}{}}
\newlabel{lemma-gysin-at-infty@cref}{{[originallemma][0][]0F9J}{[1][0][]source}{}{}{}}
\makeatother
\crefname{originallemma}{original lemma}{original lemmas}
\Crefname{originallemma}{Original lemma}{Original lemmas}
\makeatletter
\newlabel{lemma-silly}{{0F9H}{source}{}{sourcepointer-4}{}}
\newlabel{lemma-silly@cref}{{[originallemma][0][]0F9H}{[1][0][]source}{}{}{}}
\makeatother
\crefname{originallemma}{original lemma}{original lemmas}
\Crefname{originallemma}{Original lemma}{Original lemmas}
\makeatletter
\newlabel{lemma-silly-compose}{{0FAK}{source}{}{sourcepointer-5}{}}
\newlabel{lemma-silly-compose@cref}{{[originallemma][0][]0FAK}{[1][0][]source}{}{}{}}
\makeatother
\crefname{originalreference}{original reference}{original references}
\Crefname{originalreference}{Original reference}{Original references}
\makeatletter
\newlabel{remark-res-push}{{0FA1}{source}{}{sourcepointer-6}{}}
\newlabel{remark-res-push@cref}{{[originalreference][0][]0FA1}{[1][0][]source}{}{}{}}
\makeatother
\crefname{originallemma}{original lemma}{original lemmas}
\Crefname{originallemma}{Original lemma}{Original lemmas}
\makeatletter
\newlabel{lemma-homomorphism}{{0FAV}{source}{}{sourcepointer-7}{}}
\newlabel{lemma-homomorphism@cref}{{[originallemma][0][]0FAV}{[1][0][]source}{}{}{}}
\makeatother
\title{When the object vanishes over T the localized P classes are the universal polynomials in the localized Chern classes}
\author{The Stacks Project Authors}
\date{stacks-project:chow.tex:local-snapshot-20261012; GNU FDL 1.2 or later}
\begin{document}
\maketitle
\noindent\textit{Editorial note.} All delivered mathematical text is the author's own, copied from the pinned Stacks Project chapter file chow.tex; no statement or proof step was rewritten, reordered or removed.\par\medskip
\noindent\textit{Editorial note.} Preamble provenance: the pinned chapter file opens with an input of the shared preamble and carries no preamble text of its own, so the project's single Stacks-wide preamble staged at sources/stacks-chow/preamble.tex is excerpted verbatim as two spans, lines 14-18 (Xy-pic setup) and lines 154-217 (the native theorem declarations and mathematical macros). Nothing was added, removed or redefined; the wrapper class and the amsmath/amssymb/amsthm/graphicx base packages are supplied by the assembler.\par\medskip
\noindent\textit{Editorial note.} The counted statement and proof contain no cross reference, no citation and no included file, so the delivered body is self contained.\par\medskip
\noindent\textit{Editorial note.} Canonical location of the counted proposition: Stacks tag 0FAX (sources/stacks-chow/pinned-original-source.tex line 10436).\par\medskip
\noindent\textit{Allowed standard prerequisites (editorial).}\begin{itemize}
\item The localized classes P'\_p and c'\_p and the comparison with the restricted classes on E. Same chapter, Lemma labelled lemma-localized-chern-pre, cited at proof lines 10451-10462, and Lemma labelled lemma-silly-compose, cited at proof lines 10458-10462. Supplies the classes and the relation over E that the lemma transports to Z.
\item The ring structure on localized classes. Same chapter, Remark labelled remark-ring-loc-classes and Remark labelled remark-res-push, cited at proof lines 10435-10438 and 10472-10478. Defines the products appearing in the polynomial identities.
\item Pushforward of the restricted classes and the commutation with C. Same chapter, Lemma labelled lemma-silly, cited at proof lines 10484-10486, and Lemma labelled lemma-homomorphism, cited at proof lines 10488-10490. Rewrites the products through c\_\{p\_i\}(Q|\_\{W\_infinity\}) so that both sides coincide.
\item The class C of the inclusion at infinity. Same chapter, Lemma labelled lemma-gysin-at-infty, invoked at proof lines 10449-10453. Provides the class through which P'\_p and c'\_p are defined.
\end{itemize}\medskip
\ifdefined\hypertarget\hypertarget{sourcepointer-1}{}\fi
\noindent\textit{Source reference.} Original lemma 0F9K: sources/stacks-chow/pinned-original-source.tex, source lines 10204--10204. Same-chapter pointer into the author environment lemma at source line 10204; invoked at proof line 10452 of the counted proof. It is retained only as a navigation pointer; every load-bearing step is inside the delivered span.\par\medskip
\ifdefined\hypertarget\hypertarget{sourcepointer-2}{}\fi
\noindent\textit{Source reference.} Original reference 0FA0: sources/stacks-chow/pinned-original-source.tex, source lines 6419--6419. Same-chapter pointer into the author environment remark at source line 6419; invoked at proof line 10466 of the counted proof. It is retained only as a navigation pointer; every load-bearing step is inside the delivered span.\par\medskip
\ifdefined\hypertarget\hypertarget{sourcepointer-3}{}\fi
\noindent\textit{Source reference.} Original lemma 0F9J: sources/stacks-chow/pinned-original-source.tex, source lines 10026--10026. Same-chapter pointer into the author environment lemma at source line 10026; invoked at proof line 10454 of the counted proof. It is retained only as a navigation pointer; every load-bearing step is inside the delivered span.\par\medskip
\ifdefined\hypertarget\hypertarget{sourcepointer-4}{}\fi
\noindent\textit{Source reference.} Original lemma 0F9H: sources/stacks-chow/pinned-original-source.tex, source lines 9695--9695. Same-chapter pointer into the author environment lemma at source line 9695; invoked at proof line 10456 of the counted proof. It is retained only as a navigation pointer; every load-bearing step is inside the delivered span.\par\medskip
\ifdefined\hypertarget\hypertarget{sourcepointer-5}{}\fi
\noindent\textit{Source reference.} Original lemma 0FAK: sources/stacks-chow/pinned-original-source.tex, source lines 9867--9867. Same-chapter pointer into the author environment lemma at source line 9867; invoked at proof line 10458 of the counted proof. It is retained only as a navigation pointer; every load-bearing step is inside the delivered span.\par\medskip
\ifdefined\hypertarget\hypertarget{sourcepointer-6}{}\fi
\noindent\textit{Source reference.} Original reference 0FA1: sources/stacks-chow/pinned-original-source.tex, source lines 6451--6451. Same-chapter pointer into the author environment remark at source line 6451; invoked at proof line 10473 of the counted proof. It is retained only as a navigation pointer; every load-bearing step is inside the delivered span.\par\medskip
\ifdefined\hypertarget\hypertarget{sourcepointer-7}{}\fi
\noindent\textit{Source reference.} Original lemma 0FAV: sources/stacks-chow/pinned-original-source.tex, source lines 10353--10353. Same-chapter pointer into the author environment lemma at source line 10353; invoked at proof line 10483 of the counted proof. It is retained only as a navigation pointer; every load-bearing step is inside the delivered span.\par\medskip
\begin{lemma}
\label{lemma-localized-chern-pre-compose}
In Lemma \ref{lemma-localized-chern-pre} assume $Q|_T$ is zero. In
$A^*(Z \to X)$ we have
\begin{align*}
P'_1(Q) & = c'_1(Q), \\
P'_2(Q) & = c'_1(Q)^2 - 2c'_2(Q), \\
P'_3(Q) & = c'_1(Q)^3 - 3c'_1(Q)c'_2(Q) + 3c'_3(Q), \\
P'_4(Q) & = c'_1(Q)^4 - 4c'_1(Q)^2c'_2(Q) +
4c'_1(Q)c'_3(Q) + 2c'_2(Q)^2 - 4c'_4(Q),
\end{align*}
and so on with multiplication as in Remark \ref{remark-ring-loc-classes}.
\end{lemma}

\begin{proof}
The statement makes sense because the zero sheaf has rank $< 1$ and
hence the classes $c'_p(Q)$ are defined for all $p \geq 1$.
In the proof of Lemma \ref{lemma-localized-chern-pre} we have constructed
the classes $P'_p(Q)$ and $c'_p(Q)$ using the bivariant class
$C \in A^0(W_\infty \to X)$ of Lemma \ref{lemma-gysin-at-infty}
and the bivariant classes
$P'_p(Q|_E)$ and $c'_p(Q|_E)$ of Lemma \ref{lemma-silly} for the restriction
$Q|_E$ of $Q$ to the inverse image $E$ of $Z$ in $W_\infty$.
Observe that by Lemma \ref{lemma-silly-compose} we have the desired
relationship between $P'_p(Q|_E)$ and $c'_p(Q|_E)$. Recall that
$$
P'_p(Q) = (E \to Z)_* \circ P'_p(Q|_E) \circ C
\quad\text{and}\quad
c'_p(Q) = (E \to Z)_* \circ c'_p(Q|_E) \circ C
$$
To finish the proof it suffices to show the multiplications defined
in Remark \ref{remark-ring-loc-classes} on the classes $a_p = c'_p(Q)$
and on the classes $b_p = c'_p(Q|_E)$ agree:
$$
a_{p_1}a_{p_2} \ldots a_{p_r} =
(E \to Z)_* \circ b_{p_1}b_{p_2} \ldots b_{p_r} \circ C
$$
Some details omitted. If $r = 1$, then this is true.
For $r > 1$ note that by Remark \ref{remark-res-push} the multiplication in
Remark \ref{remark-ring-loc-classes} proceeds
by inserting $(Z \to X)_*$, resp.\ $(E \to W_\infty)_*$ in between
the factors of the product
$a_{p_1}a_{p_2} \ldots a_{p_r}$, resp.\ $b_{p_1}b_{p_2} \ldots b_{p_r}$
and taking compositions as bivariant classes.
Now by Lemma \ref{lemma-silly} we have
$$
(E \to W_\infty)_* \circ b_{p_i} = c_{p_i}(Q|_{W_\infty})
$$
and by Lemma \ref{lemma-homomorphism} we have
$$
C \circ (Z \to X)_* \circ a_{p_i} = c_{p_i}(Q|_{W_\infty}) \circ C
$$
for $i = 2, \ldots, r$. A calculation
shows that the left and right hand side of the desired
equality both simplify to
$$
(E \to Z)_* \circ c'_{p_1}(Q|_E) \circ
c_{p_2}(Q|_{W_\infty}) \circ \ldots \circ
c_{p_r}(Q|_{W_\infty}) \circ C
$$
and the proof is complete.
\end{proof}

\end{document}
